% =====================================================================
% ANLAGENTEIL -- Nummerierung und Bezeichnung folgen den Abschnitten
% (Gibt es zu einem Abschnitt keine Anlage, entfaellt die Nummer.)
% =====================================================================

% ---------------------------------------------------------------------
\setzeTeil{Anlage 1}
\section{Allgemeines}\label{sec:anlagenteil1}

Anlagen zu Abschnitt~1: Leistungsversprechen und Selbstbehalte des
Tarifs.

\setzeTeil{Anlage 1.1}
\subsection{Allgemeine Darstellung des Tarifs}\label{sec:anlage1}

Maßgeblich sind die AVB (Teil~I MB/KK~2009, Teil~II Tarifbedingungen
VitalPlus, Stand 01/2019 -- letzte bedingungsseitige Änderung mit
Änderungsnr.~11; die AVB liegen dieser TB als gesondertes Dokument
bei). Kurzübersicht des Leistungsversprechens:

\begin{center}\small
\begin{tabular}{@{}lp{9.5cm}@{}}
\toprule
Leistungsbereich & Erstattung \\
\midrule
ambulante Heilbehandlung & 100\,\% nach GOÄ bis zu den Höchstsätzen; absoluter Selbstbehalt 100~EUR je Kalenderjahr (Kinder/Jugendliche: 0~EUR) \\
stationäre Heilbehandlung & 100\,\% allgemeine Krankenhausleistungen, Unterbringung im Mehrbettzimmer \\
Zahnbehandlung & 100\,\% \\
Zahnersatz und Kieferorthopädie & 75\,\% (Erstattungsklasse 70--80\,\%) \\
\bottomrule
\end{tabular}
\end{center}

\textbf{Höhe und Entwicklung der Selbstbehalte:} Der absolute
Selbstbehalt im ambulanten Bereich beträgt seit Tarifeinführung
(01.01.1998) unverändert 100~EUR je Kalenderjahr für Erwachsene und
0~EUR für Kinder und Jugendliche; alle Beobachtungseinheiten führen
dieselbe Selbstbehaltsstufe.

% ---------------------------------------------------------------------
\setzeTeil{Anlage 2}
\section{Rechnungsgrundlagen}\label{sec:anlagenteil2}

Anlagen zu Abschnitt~2: Rechnungszinssätze, Ausscheideordnung,
Kopfschäden und Zuschläge.

\setzeTeil{Anlage 2.1}
\subsection{Rechnungszinssatz}\label{sec:anlage21}

Rechnungszinssätze $i$ je Beobachtungseinheit (jeweils unterhalb des
AUZ von 2{,}8\,\%, Nachweis in der Rahmen-TB RTB-AUZ 2025-1):

\begin{center}
\begin{tabular}{ccc}
\toprule
Kinder/Jugendliche & Männer & Frauen \\
\midrule
$\ZinsK\,\%$ & $\ZinsM\,\%$ & $\ZinsF\,\%$ \\
\bottomrule
\end{tabular}
\end{center}

Historie: Männer 3{,}0\,\% $\to$ 2{,}7\,\% (Änderung Nr.\,14); Frauen
2{,}3\,\% $\to$ 2{,}0\,\% (Änderung Nr.\,15); Kinder/Jugendliche
unverändert seit Änderung Nr.\,12.

% ---------------------------------------------------------------------
\setzeTeil{Anlage 2.2}
\subsection{Ausscheideordnung}\label{sec:anlage221}

Rechnungsmäßige Sterbewahrscheinlichkeiten $q_x$ (PKV-Sterbetafel 2025,
Fortschreibung gemäß Gleichung~\eqref{eq:qfort}),
Stornowahrscheinlichkeiten $w_x$ (unternehmenseigene Beobachtung,
Abschnitt~\ref{sec:storno}), Ausscheidewahrscheinlichkeiten und
diskontierte Übergangswahrscheinlichkeiten
$DQ_x = v\,(1-(q+w)_x)$. Endalter: 110.

\textbf{Nachweis der Stornowahrscheinlichkeiten:} Datenbasis der
unternehmenseigenen Beobachtung 2021--2023 und Vergleich mit den
bisher angesetzten Werten (Änderungsnr.~16); alle Angaben fiktiv:

\begin{center}\small
\begin{tabular}{@{}lrccc@{}}
\toprule
Altersband & Bestand (Personenjahre) & $\bar{w}$ beobachtet & $\bar{w}$ bisher & $\bar{w}$ neu \\
\midrule
21--25 & 14\,830 & 0{,}0994 & 0{,}1102 & 0{,}1054 \\
26--35 & 41\,270 & 0{,}0801 & 0{,}0898 & 0{,}0846 \\
36--50 & 68\,930 & 0{,}0392 & 0{,}0391 & 0{,}0391 \\
51--65 & 55\,410 & 0{,}0181 & 0{,}0184 & 0{,}0184 \\
66--110 & 21\,060 & 0{,}0021 & 0{,}0022 & 0{,}0022 \\
\bottomrule
\end{tabular}
\end{center}

Die rechnungsmäßigen Werte $\bar{w}$ neu liegen oberhalb der
beobachteten Rohwerte (Sicherheitsmarge, da eine Überschätzung des
Stornos die Vererbung überzeichnen würde, vgl.
Abschnitt~\ref{sec:tarifar}); in den Altern 21--35 wurden sie
gegenüber dem bisherigen Ansatz um durchschnittlich
0{,}5~Prozentpunkte gesenkt.

% Automatisch erzeugt von kalkulation/pipeline.py
\begin{small}
\begin{longtable}{rrrrr}
\caption{Ausscheideordnung Männer: $q_x$, $w_x$, $(q+w)_x$ und $DQ_x$ (Rechnungszins 2,7\,\%)}\label{tab:aoM}\\
\toprule
$x$ & $q_x$ & $w_x$ & $(q+w)_x$ & $DQ_x$ \\
\midrule
\endfirsthead
\caption[]{Ausscheideordnung Männer: $q_x$, $w_x$, $(q+w)_x$ und $DQ_x$ (Rechnungszins 2,7\,\%) (Fortsetzung)}\\
\toprule
$x$ & $q_x$ & $w_x$ & $(q+w)_x$ & $DQ_x$ \\
\midrule
\endhead
\bottomrule
\endfoot
21 & 0,000382 & 0,105000 & 0,105382 & 0,871098 \\
22 & 0,000402 & 0,115000 & 0,115402 & 0,861342 \\
23 & 0,000414 & 0,125000 & 0,125414 & 0,851593 \\
24 & 0,000398 & 0,117000 & 0,117398 & 0,859398 \\
25 & 0,000380 & 0,109000 & 0,109380 & 0,867205 \\
26 & 0,000359 & 0,101000 & 0,101359 & 0,875016 \\
27 & 0,000333 & 0,093000 & 0,093333 & 0,882831 \\
28 & 0,000311 & 0,085000 & 0,085311 & 0,890642 \\
29 & 0,000296 & 0,078000 & 0,078296 & 0,897472 \\
30 & 0,000288 & 0,071000 & 0,071288 & 0,904296 \\
31 & 0,000286 & 0,064000 & 0,064286 & 0,911114 \\
32 & 0,000290 & 0,057000 & 0,057290 & 0,917926 \\
33 & 0,000293 & 0,050000 & 0,050293 & 0,924739 \\
34 & 0,000303 & 0,048000 & 0,048303 & 0,926677 \\
35 & 0,000318 & 0,046000 & 0,046318 & 0,928610 \\
36 & 0,000330 & 0,044000 & 0,044330 & 0,930545 \\
37 & 0,000348 & 0,042000 & 0,042348 & 0,932475 \\
38 & 0,000373 & 0,040000 & 0,040373 & 0,934398 \\
39 & 0,000404 & 0,039000 & 0,039404 & 0,935342 \\
40 & 0,000441 & 0,038000 & 0,038441 & 0,936279 \\
41 & 0,000485 & 0,037000 & 0,037485 & 0,937210 \\
42 & 0,000536 & 0,036000 & 0,036536 & 0,938134 \\
43 & 0,000592 & 0,035000 & 0,035592 & 0,939054 \\
44 & 0,000636 & 0,035000 & 0,035636 & 0,939011 \\
45 & 0,000687 & 0,035000 & 0,035687 & 0,938961 \\
46 & 0,000757 & 0,035000 & 0,035757 & 0,938893 \\
47 & 0,000850 & 0,035000 & 0,035850 & 0,938802 \\
48 & 0,000968 & 0,035000 & 0,035968 & 0,938687 \\
49 & 0,001109 & 0,035000 & 0,036109 & 0,938550 \\
50 & 0,001271 & 0,035000 & 0,036271 & 0,938392 \\
51 & 0,001448 & 0,035000 & 0,036448 & 0,938220 \\
52 & 0,001634 & 0,035000 & 0,036634 & 0,938039 \\
53 & 0,001840 & 0,035000 & 0,036840 & 0,937838 \\
54 & 0,002073 & 0,028600 & 0,030673 & 0,943843 \\
55 & 0,002338 & 0,022200 & 0,024538 & 0,949817 \\
56 & 0,002644 & 0,015800 & 0,018444 & 0,955751 \\
57 & 0,002994 & 0,009400 & 0,012394 & 0,961642 \\
58 & 0,003394 & 0,003000 & 0,006394 & 0,967484 \\
59 & 0,003850 & 0,002800 & 0,006650 & 0,967235 \\
60 & 0,004234 & 0,002600 & 0,006834 & 0,967056 \\
61 & 0,004616 & 0,002400 & 0,007016 & 0,966878 \\
62 & 0,005023 & 0,002200 & 0,007223 & 0,966677 \\
63 & 0,005416 & 0,002000 & 0,007416 & 0,966489 \\
64 & 0,005877 & 0,001800 & 0,007677 & 0,966235 \\
65 & 0,006425 & 0,001600 & 0,008025 & 0,965896 \\
66 & 0,007037 & 0,001400 & 0,008437 & 0,965495 \\
67 & 0,007701 & 0,001200 & 0,008901 & 0,965043 \\
68 & 0,008431 & 0,001000 & 0,009431 & 0,964527 \\
69 & 0,009197 & 0,001000 & 0,010197 & 0,963781 \\
70 & 0,010079 & 0,001000 & 0,011079 & 0,962922 \\
71 & 0,011107 & 0,001000 & 0,012107 & 0,961921 \\
72 & 0,012282 & 0,001000 & 0,013282 & 0,960777 \\
73 & 0,013674 & 0,001000 & 0,014674 & 0,959422 \\
74 & 0,015323 & 0,001000 & 0,016323 & 0,957816 \\
75 & 0,017251 & 0,001000 & 0,018251 & 0,955939 \\
76 & 0,019376 & 0,001000 & 0,020376 & 0,953870 \\
77 & 0,021526 & 0,001000 & 0,022526 & 0,951776 \\
78 & 0,024097 & 0,001000 & 0,025097 & 0,949273 \\
79 & 0,027194 & 0,001000 & 0,028194 & 0,946257 \\
80 & 0,030930 & 0,001000 & 0,031930 & 0,942619 \\
81 & 0,035428 & 0,001000 & 0,036428 & 0,938240 \\
82 & 0,040788 & 0,001000 & 0,041788 & 0,933020 \\
83 & 0,047012 & 0,001000 & 0,048012 & 0,926960 \\
84 & 0,054397 & 0,001000 & 0,055397 & 0,919769 \\
85 & 0,063080 & 0,001000 & 0,064080 & 0,911315 \\
86 & 0,073180 & 0,001000 & 0,074180 & 0,901480 \\
87 & 0,084803 & 0,001000 & 0,085803 & 0,890163 \\
88 & 0,098028 & 0,001000 & 0,099028 & 0,877285 \\
89 & 0,112913 & 0,001000 & 0,113913 & 0,862792 \\
90 & 0,129465 & 0,001000 & 0,130465 & 0,846675 \\
91 & 0,147647 & 0,001000 & 0,148647 & 0,828971 \\
92 & 0,167363 & 0,001000 & 0,168363 & 0,809773 \\
93 & 0,188107 & 0,001000 & 0,189107 & 0,789574 \\
94 & 0,209873 & 0,001000 & 0,210873 & 0,768381 \\
95 & 0,232032 & 0,001000 & 0,233032 & 0,746804 \\
96 & 0,243250 & 0,001000 & 0,244250 & 0,735881 \\
97 & 0,251290 & 0,001000 & 0,252290 & 0,728053 \\
98 & 0,258960 & 0,001000 & 0,259960 & 0,720584 \\
99 & 0,266210 & 0,001000 & 0,267210 & 0,713525 \\
100 & 0,273020 & 0,001000 & 0,274020 & 0,706894 \\
101 & 0,279340 & 0,001000 & 0,280340 & 0,700740 \\
102 & 0,285140 & 0,001000 & 0,286140 & 0,695093 \\
103 & 0,374498 & 0,001000 & 0,375498 & 0,608084 \\
104 & 0,463855 & 0,001000 & 0,464855 & 0,521076 \\
105 & 0,553213 & 0,001000 & 0,554213 & 0,434067 \\
106 & 0,642570 & 0,001000 & 0,643570 & 0,347059 \\
107 & 0,731928 & 0,001000 & 0,732928 & 0,260051 \\
108 & 0,821285 & 0,001000 & 0,822285 & 0,173043 \\
109 & 0,910643 & 0,001000 & 0,911643 & 0,086034 \\
110 & 1,000000 & 0,000000 & 1,000000 & 0,000000 \\
\end{longtable}
\end{small}

\clearpage
% Automatisch erzeugt von kalkulation/pipeline.py
\begin{small}
\begin{longtable}{rrrrr}
\caption{Ausscheideordnung Frauen: $q_x$, $w_x$, $(q+w)_x$ und $DQ_x$ (Rechnungszins 2,0\,\%)}\label{tab:aoF}\\
\toprule
$x$ & $q_x$ & $w_x$ & $(q+w)_x$ & $DQ_x$ \\
\midrule
\endfirsthead
\caption[]{Ausscheideordnung Frauen: $q_x$, $w_x$, $(q+w)_x$ und $DQ_x$ (Rechnungszins 2,0\,\%) (Fortsetzung)}\\
\toprule
$x$ & $q_x$ & $w_x$ & $(q+w)_x$ & $DQ_x$ \\
\midrule
\endhead
\bottomrule
\endfoot
21 & 0,000203 & 0,105000 & 0,105203 & 0,877252 \\
22 & 0,000202 & 0,115000 & 0,115202 & 0,867449 \\
23 & 0,000200 & 0,125000 & 0,125200 & 0,857647 \\
24 & 0,000197 & 0,117000 & 0,117197 & 0,865493 \\
25 & 0,000195 & 0,109000 & 0,109195 & 0,873338 \\
26 & 0,000196 & 0,101000 & 0,101196 & 0,881180 \\
27 & 0,000194 & 0,093000 & 0,093194 & 0,889025 \\
28 & 0,000184 & 0,085000 & 0,085184 & 0,896878 \\
29 & 0,000178 & 0,078000 & 0,078178 & 0,903747 \\
30 & 0,000176 & 0,071000 & 0,071176 & 0,910612 \\
31 & 0,000177 & 0,064000 & 0,064177 & 0,917474 \\
32 & 0,000181 & 0,057000 & 0,057181 & 0,924332 \\
33 & 0,000187 & 0,050000 & 0,050187 & 0,931189 \\
34 & 0,000194 & 0,048000 & 0,048194 & 0,933143 \\
35 & 0,000203 & 0,046000 & 0,046203 & 0,935095 \\
36 & 0,000218 & 0,044000 & 0,044218 & 0,937041 \\
37 & 0,000237 & 0,042000 & 0,042237 & 0,938983 \\
38 & 0,000262 & 0,040000 & 0,040262 & 0,940920 \\
39 & 0,000291 & 0,039000 & 0,039291 & 0,941872 \\
40 & 0,000321 & 0,038000 & 0,038321 & 0,942823 \\
41 & 0,000357 & 0,037000 & 0,037357 & 0,943768 \\
42 & 0,000399 & 0,036000 & 0,036399 & 0,944707 \\
43 & 0,000444 & 0,035000 & 0,035444 & 0,945643 \\
44 & 0,000486 & 0,035000 & 0,035486 & 0,945602 \\
45 & 0,000532 & 0,035000 & 0,035532 & 0,945557 \\
46 & 0,000586 & 0,035000 & 0,035586 & 0,945504 \\
47 & 0,000648 & 0,035000 & 0,035648 & 0,945443 \\
48 & 0,000719 & 0,035000 & 0,035719 & 0,945374 \\
49 & 0,000799 & 0,035000 & 0,035799 & 0,945295 \\
50 & 0,000887 & 0,035000 & 0,035887 & 0,945209 \\
51 & 0,000983 & 0,035000 & 0,035983 & 0,945115 \\
52 & 0,001083 & 0,035000 & 0,036083 & 0,945017 \\
53 & 0,001193 & 0,035000 & 0,036193 & 0,944909 \\
54 & 0,001313 & 0,028600 & 0,029913 & 0,951066 \\
55 & 0,001449 & 0,022200 & 0,023649 & 0,957207 \\
56 & 0,001604 & 0,015800 & 0,017404 & 0,963329 \\
57 & 0,001782 & 0,009400 & 0,011182 & 0,969429 \\
58 & 0,001987 & 0,003000 & 0,004987 & 0,975503 \\
59 & 0,002222 & 0,002800 & 0,005022 & 0,975469 \\
60 & 0,002488 & 0,002600 & 0,005088 & 0,975404 \\
61 & 0,002785 & 0,002400 & 0,005185 & 0,975309 \\
62 & 0,003118 & 0,002200 & 0,005318 & 0,975178 \\
63 & 0,003483 & 0,002000 & 0,005483 & 0,975017 \\
64 & 0,003795 & 0,001800 & 0,005595 & 0,974907 \\
65 & 0,004092 & 0,001600 & 0,005692 & 0,974812 \\
66 & 0,004406 & 0,001400 & 0,005806 & 0,974700 \\
67 & 0,004751 & 0,001200 & 0,005951 & 0,974558 \\
68 & 0,005149 & 0,001000 & 0,006149 & 0,974364 \\
69 & 0,005618 & 0,001000 & 0,006618 & 0,973904 \\
70 & 0,006131 & 0,001000 & 0,007131 & 0,973401 \\
71 & 0,006710 & 0,001000 & 0,007710 & 0,972833 \\
72 & 0,007376 & 0,001000 & 0,008376 & 0,972180 \\
73 & 0,008175 & 0,001000 & 0,009175 & 0,971397 \\
74 & 0,009154 & 0,001000 & 0,010154 & 0,970437 \\
75 & 0,010351 & 0,001000 & 0,011351 & 0,969264 \\
76 & 0,011808 & 0,001000 & 0,012808 & 0,967835 \\
77 & 0,013297 & 0,001000 & 0,014297 & 0,966375 \\
78 & 0,014989 & 0,001000 & 0,015989 & 0,964717 \\
79 & 0,017113 & 0,001000 & 0,018113 & 0,962634 \\
80 & 0,019770 & 0,001000 & 0,020770 & 0,960029 \\
81 & 0,023054 & 0,001000 & 0,024054 & 0,956810 \\
82 & 0,026977 & 0,001000 & 0,027977 & 0,952964 \\
83 & 0,031755 & 0,001000 & 0,032755 & 0,948279 \\
84 & 0,037512 & 0,001000 & 0,038512 & 0,942635 \\
85 & 0,044371 & 0,001000 & 0,045371 & 0,935911 \\
86 & 0,052454 & 0,001000 & 0,053454 & 0,927986 \\
87 & 0,061881 & 0,001000 & 0,062881 & 0,918744 \\
88 & 0,072769 & 0,001000 & 0,073769 & 0,908070 \\
89 & 0,085225 & 0,001000 & 0,086225 & 0,895858 \\
90 & 0,099336 & 0,001000 & 0,100336 & 0,882024 \\
91 & 0,115161 & 0,001000 & 0,116161 & 0,866509 \\
92 & 0,132712 & 0,001000 & 0,133712 & 0,849302 \\
93 & 0,151809 & 0,001000 & 0,152809 & 0,830579 \\
94 & 0,169682 & 0,001000 & 0,170682 & 0,813057 \\
95 & 0,188536 & 0,001000 & 0,189536 & 0,794573 \\
96 & 0,208205 & 0,001000 & 0,209205 & 0,775289 \\
97 & 0,228513 & 0,001000 & 0,229513 & 0,755379 \\
98 & 0,249230 & 0,001000 & 0,250230 & 0,735069 \\
99 & 0,258450 & 0,001000 & 0,259450 & 0,726029 \\
100 & 0,262980 & 0,001000 & 0,263980 & 0,721588 \\
101 & 0,268110 & 0,001000 & 0,269110 & 0,716559 \\
102 & 0,272880 & 0,001000 & 0,273880 & 0,711882 \\
103 & 0,363770 & 0,001000 & 0,364770 & 0,622775 \\
104 & 0,454660 & 0,001000 & 0,455660 & 0,533667 \\
105 & 0,545550 & 0,001000 & 0,546550 & 0,444559 \\
106 & 0,636440 & 0,001000 & 0,637440 & 0,355451 \\
107 & 0,727330 & 0,001000 & 0,728330 & 0,266343 \\
108 & 0,818220 & 0,001000 & 0,819220 & 0,177235 \\
109 & 0,909110 & 0,001000 & 0,910110 & 0,088127 \\
110 & 1,000000 & 0,000000 & 1,000000 & 0,000000 \\
\end{longtable}
\end{small}


% ---------------------------------------------------------------------
\setzeTeil{Anlage 2.3}
\subsection{Kopfschaden}\label{sec:anlagenteil23}

Die Unteranlagen dokumentieren die Herleitung der rechnungsmäßigen
Kopfschäden gemäß Abschnitt~\ref{sec:kopfschaden} aus den öffentlichen
Stützdaten (BaFin-Kopfschadenstatistik, Beobachtungsjahr 2022).

\setzeTeil{Anlage 2.3.2}
\setcounter{subsubsection}{1}
\subsubsection{Umlage der Schwangerschafts- und Mutterschaftskosten}\label{sec:anlage232}\markright{Anlage~2.3.2}

Beobachtete Teilkopfschäden je Leistungsbereich nach Umlage der
S-Kosten gemäß Gleichungen~\eqref{eq:sumlage}--\eqref{eq:ksumme}
(Quelle: BaFin-Kopfschadenstatistik, Beobachtungsjahr 2022; Tafeln
gemäß Abschnitt~\ref{sec:kopfschaden-grundsaetze}; Tafelbereich
Alter 0--100, Fortschreibung der Alter 101--110 gemäß
Gleichung~\eqref{eq:kfort}). Zur Nachvollziehbarkeit der Umlage sind
je Alter die Summe der S-Kosten $\sum K^{S}_x$ (Frauentafel) und die
Summe der Umlagen $\sum u_x$ über die vier Leistungsbereiche
angegeben; die Rohwerte vor Umlage sowie die in
Gleichung~\eqref{eq:sumlage} verwendeten Bestände $n^{M}_x$, $n^{W}_x$
(Gewichte, auch für die Gruppenmittel der Kinder und Jugendlichen)
sind die veröffentlichten Werte der genannten BaFin-Tafeln (Spalten
KOPFSCHADEN und BESTAND) und werden hier nicht abgedruckt. Die
Summenspalte ist Grundlage der Profile und Grundkopfschäden in
Anlage~\ref{sec:anlage233}.

Die für die Fortschreibung der Alter 101--110
(Gleichung~\eqref{eq:kfort}) maßgebliche mittlere jährliche
Veränderungsrate der Gesamtkopfschäden der Alter 95 bis 100 beträgt
\WachstumM\,\% p.\,a. (Männer) bzw. \WachstumF\,\% p.\,a. (Frauen).

% Automatisch erzeugt von kalkulation/pipeline.py
\begin{small}
\begin{longtable}{rrrrrrr}
\caption{Herleitung der Kopfschäden Männer: beobachtete Teilkopfschäden je Leistungsbereich nach S-Kosten-Umlage (EUR p.\,a., Beobachtungsjahr 2022, Tafelbereich Alter 0--100)}\label{tab:herleitungM}\\
\toprule
$x$ & ambulant & stationär & Zahn-Beh. & Zahn-Ers. & $\sum u_x$ & Summe $K_x^{\mathrm{beob}}$ \\
\midrule
\endfirsthead
\caption[]{Herleitung der Kopfschäden Männer: beobachtete Teilkopfschäden je Leistungsbereich nach S-Kosten-Umlage (EUR p.\,a., Beobachtungsjahr 2022, Tafelbereich Alter 0--100) (Fortsetzung)}\\
\toprule
$x$ & ambulant & stationär & Zahn-Beh. & Zahn-Ers. & $\sum u_x$ & Summe $K_x^{\mathrm{beob}}$ \\
\midrule
\endhead
\bottomrule
\endfoot
0 & 1247,3906 & 2572,7315 & 0,0000 & 0,6486 & 0,0000 & 3820,7707 \\
1 & 1168,0867 & 575,3253 & 9,0216 & 0,0000 & 0,0000 & 1752,4336 \\
2 & 1091,4336 & 423,2235 & 22,8722 & 0,0000 & 0,0000 & 1537,5293 \\
3 & 1021,8810 & 310,3731 & 41,2701 & 0,0000 & 0,0000 & 1373,5242 \\
4 & 962,4758 & 238,3399 & 63,5802 & 0,2234 & 0,0000 & 1264,6193 \\
5 & 914,6624 & 190,7164 & 88,3674 & 8,2075 & 0,0000 & 1201,9537 \\
6 & 878,7477 & 172,2433 & 113,2455 & 27,4322 & 0,0000 & 1191,6687 \\
7 & 854,7810 & 167,6676 & 135,2944 & 64,0739 & 0,0000 & 1221,8169 \\
8 & 842,7021 & 163,4999 & 152,1431 & 119,5400 & 0,0000 & 1277,8851 \\
9 & 841,9588 & 154,7004 & 163,2927 & 190,9412 & 0,0000 & 1350,8931 \\
10 & 851,2848 & 142,8012 & 170,5237 & 285,3796 & 0,0000 & 1449,9893 \\
11 & 869,2889 & 138,1069 & 176,8033 & 426,9694 & 0,0000 & 1611,1685 \\
12 & 894,2251 & 147,9356 & 184,5381 & 616,2166 & 0,0000 & 1842,9154 \\
13 & 924,0903 & 171,7786 & 194,2280 & 796,8824 & 0,0000 & 2086,9793 \\
14 & 956,7338 & 201,4442 & 204,5037 & 882,4174 & 0,0000 & 2245,0991 \\
15 & 990,3058 & 240,9020 & 213,1166 & 826,7588 & 0,0138 & 2271,0832 \\
16 & 1022,6762 & 288,0960 & 217,9624 & 663,1041 & 0,0183 & 2191,8387 \\
17 & 1052,2085 & 338,7932 & 217,8996 & 466,9707 & 0,4278 & 2075,8719 \\
18 & 1077,7046 & 398,3393 & 213,1098 & 303,0506 & 0,3925 & 1992,2044 \\
19 & 1099,1843 & 462,7284 & 204,9192 & 169,7186 & 0,6862 & 1936,5505 \\
20 & 1116,6670 & 505,2775 & 195,0884 & 148,4016 & 1,7548 & 1965,4346 \\
21 & 1130,1793 & 528,3758 & 185,2659 & 128,4469 & 4,1329 & 1972,2680 \\
22 & 1139,8733 & 519,9020 & 176,6961 & 111,1288 & 7,1720 & 1947,6002 \\
23 & 1149,0854 & 496,6713 & 170,0950 & 97,2980 & 13,0815 & 1913,1498 \\
24 & 1159,8671 & 455,8310 & 165,6734 & 87,1078 & 23,9583 & 1868,4794 \\
25 & 1174,2582 & 395,0297 & 163,2403 & 81,0986 & 44,3334 & 1813,6268 \\
26 & 1193,2913 & 327,4313 & 162,3260 & 79,1049 & 78,2425 & 1762,1535 \\
27 & 1212,4936 & 334,0974 & 162,3212 & 80,1203 & 110,9447 & 1789,0324 \\
28 & 1228,1407 & 332,2744 & 162,7273 & 82,8040 & 137,7244 & 1805,9463 \\
29 & 1246,7637 & 332,8099 & 163,3049 & 86,0266 & 170,7295 & 1828,9052 \\
30 & 1259,1574 & 330,3016 & 164,0313 & 90,0121 & 194,3098 & 1843,5024 \\
31 & 1277,5714 & 332,8186 & 165,0121 & 95,2022 & 227,7105 & 1870,6044 \\
32 & 1281,1613 & 329,1348 & 166,3450 & 101,5555 & 237,8214 & 1878,1966 \\
33 & 1290,7867 & 328,6899 & 168,1012 & 107,6920 & 253,8867 & 1895,2699 \\
34 & 1283,6601 & 328,1561 & 170,1963 & 114,2755 & 248,7523 & 1896,2882 \\
35 & 1280,2276 & 326,7996 & 172,4350 & 123,4019 & 241,7082 & 1902,8641 \\
36 & 1274,0236 & 322,2458 & 174,6323 & 134,6973 & 223,6990 & 1905,5989 \\
37 & 1267,0816 & 318,8959 & 176,6880 & 146,7558 & 201,2536 & 1909,4213 \\
38 & 1254,5833 & 311,5161 & 178,5677 & 143,2718 & 164,0130 & 1887,9388 \\
39 & 1254,3606 & 309,8484 & 180,3492 & 151,9482 & 139,2546 & 1896,5064 \\
40 & 1252,1798 & 306,0499 & 182,2281 & 160,5492 & 104,4367 & 1901,0070 \\
41 & 1260,9782 & 305,5012 & 184,4405 & 169,1900 & 77,7675 & 1920,1099 \\
42 & 1273,5546 & 310,3628 & 187,2233 & 178,1973 & 53,2323 & 1949,3379 \\
43 & 1295,9734 & 319,0011 & 190,8851 & 187,5817 & 35,3346 & 1993,4413 \\
44 & 1326,6103 & 332,0106 & 195,4878 & 197,2182 & 22,9330 & 2051,3269 \\
45 & 1364,1362 & 348,2745 & 200,6085 & 207,1034 & 13,5757 & 2120,1225 \\
46 & 1408,3673 & 367,3799 & 205,6178 & 217,2946 & 7,9051 & 2198,6595 \\
47 & 1457,7489 & 388,6159 & 210,0158 & 227,9653 & 4,3485 & 2284,3459 \\
48 & 1511,8816 & 412,8421 & 213,6704 & 239,3308 & 2,5789 & 2377,7249 \\
49 & 1569,8078 & 440,7042 & 216,8937 & 251,2210 & 1,7519 & 2478,6268 \\
50 & 1631,3665 & 472,1528 & 220,2335 & 263,5619 & 1,3834 & 2587,3148 \\
51 & 1696,9041 & 506,1029 & 224,0838 & 276,3771 & 1,2428 & 2703,4679 \\
52 & 1767,1183 & 543,0390 & 228,4138 & 289,8037 & 1,1271 & 2828,3748 \\
53 & 1843,4692 & 584,2576 & 232,6006 & 303,6779 & 1,1103 & 2964,0053 \\
54 & 1926,8215 & 629,4369 & 236,0704 & 317,4801 & 0,7681 & 3109,8089 \\
55 & 2017,4569 & 678,5116 & 238,5374 & 331,3156 & 0,6293 & 3265,8215 \\
56 & 2115,8493 & 731,4033 & 240,3768 & 345,7069 & 0,4404 & 3433,3363 \\
57 & 2222,1668 & 788,1535 & 242,2663 & 361,2512 & 0,2594 & 3613,8378 \\
58 & 2336,4471 & 848,7570 & 244,7005 & 378,1426 & 0,1311 & 3808,0472 \\
59 & 2459,7008 & 913,4675 & 247,8016 & 395,9382 & 0,0477 & 4016,9081 \\
60 & 2592,9399 & 983,8830 & 251,4807 & 414,2343 & 0,0012 & 4242,5379 \\
61 & 2736,1464 & 1060,9738 & 255,7980 & 433,1075 & 0,0000 & 4486,0257 \\
62 & 2888,2868 & 1146,2847 & 260,9094 & 452,4999 & 0,0000 & 4747,9808 \\
63 & 3048,4742 & 1240,5562 & 266,6691 & 472,0087 & 0,0000 & 5027,7082 \\
64 & 3216,6409 & 1343,1391 & 272,3593 & 490,5143 & 0,0000 & 5322,6536 \\
65 & 3392,0977 & 1452,8684 & 277,0747 & 506,9726 & 0,0000 & 5629,0134 \\
66 & 3573,6064 & 1567,6589 & 280,4103 & 520,5919 & 0,0000 & 5942,2675 \\
67 & 3760,2133 & 1683,6382 & 282,4682 & 531,4914 & 0,0000 & 6257,8111 \\
68 & 3951,8219 & 1800,3866 & 283,6474 & 540,1314 & 0,0000 & 6575,9873 \\
69 & 4148,7364 & 1919,1607 & 284,5967 & 547,0374 & 0,0000 & 6899,5312 \\
70 & 4351,3833 & 2040,8627 & 285,8953 & 553,0045 & 0,0000 & 7231,1458 \\
71 & 4558,5788 & 2165,1146 & 287,9994 & 558,5081 & 0,0000 & 7570,2009 \\
72 & 4768,2429 & 2290,5976 & 290,9218 & 563,9744 & 0,0000 & 7913,7367 \\
73 & 4978,5710 & 2416,3344 & 294,1883 & 569,3536 & 0,0000 & 8258,4473 \\
74 & 5187,2351 & 2541,3905 & 297,1913 & 574,6698 & 0,0000 & 8600,4867 \\
75 & 5391,1981 & 2664,8499 & 299,5888 & 579,7734 & 0,0000 & 8935,4102 \\
76 & 5611,9538 & 2787,1783 & 301,3672 & 584,4562 & 0,0000 & 9284,9555 \\
77 & 5807,2513 & 2909,7630 & 302,5551 & 588,0230 & 0,0000 & 9607,5924 \\
78 & 5990,0621 & 3032,6284 & 303,1193 & 589,8359 & 0,0000 & 9915,6457 \\
79 & 6158,2320 & 3155,0014 & 302,9308 & 589,2167 & 0,0000 & 10205,3809 \\
80 & 6308,8706 & 3275,4683 & 301,8990 & 585,9534 & 0,0000 & 10472,1913 \\
81 & 6440,7702 & 3406,3013 & 300,0877 & 580,0882 & 0,0000 & 10727,2474 \\
82 & 6553,1284 & 3538,4437 & 297,4936 & 571,9759 & 0,0000 & 10961,0416 \\
83 & 6645,6802 & 3663,8152 & 293,8877 & 562,0261 & 0,0000 & 11165,4092 \\
84 & 6718,7260 & 3776,4573 & 288,7947 & 550,4168 & 0,0000 & 11334,3948 \\
85 & 6772,3917 & 3872,0362 & 281,7624 & 537,0281 & 0,0000 & 11463,2184 \\
86 & 6807,7333 & 3948,6892 & 272,5706 & 521,8081 & 0,0000 & 11550,8012 \\
87 & 6827,2285 & 4006,4191 & 261,3432 & 504,8191 & 0,0000 & 11599,8099 \\
88 & 6833,6111 & 4045,3311 & 248,4725 & 486,2589 & 0,0000 & 11613,6736 \\
89 & 6833,6111 & 4065,5466 & 234,4448 & 466,4267 & 0,0000 & 11600,0292 \\
90 & 6833,6111 & 4067,4705 & 219,7415 & 445,6200 & 0,0000 & 11566,4431 \\
91 & 6833,6111 & 4067,4705 & 204,7774 & 424,0388 & 0,0000 & 11529,8978 \\
92 & 6833,6111 & 4067,4705 & 189,8743 & 401,9122 & 0,0000 & 11492,8681 \\
93 & 6833,6111 & 4067,4705 & 175,2514 & 379,4435 & 0,0000 & 11455,7765 \\
94 & 6833,6111 & 4067,4705 & 161,0166 & 356,7565 & 0,0000 & 11418,8547 \\
95 & 6833,6111 & 4067,4705 & 147,2359 & 333,9403 & 0,0000 & 11382,2578 \\
96 & 6833,6111 & 4067,4705 & 133,9270 & 311,0620 & 0,0000 & 11346,0706 \\
97 & 6833,6111 & 4067,4705 & 121,0904 & 288,1598 & 0,0000 & 11310,3318 \\
98 & 6833,6111 & 4067,4705 & 108,7279 & 265,2449 & 0,0000 & 11275,0544 \\
99 & 6833,6111 & 4067,4705 & 96,8398 & 242,3277 & 0,0000 & 11240,2491 \\
100 & 6833,6111 & 4067,4705 & 85,4276 & 219,4104 & 0,0000 & 11205,9196 \\
\end{longtable}
\end{small}

\clearpage
% Automatisch erzeugt von kalkulation/pipeline.py
\begin{small}
\begin{longtable}{rrrrrrrr}
\caption{Herleitung der Kopfschäden Frauen: beobachtete Teilkopfschäden je Leistungsbereich nach S-Kosten-Umlage (EUR p.\,a., Beobachtungsjahr 2022, Tafelbereich Alter 0--100)}\label{tab:herleitungF}\\
\toprule
$x$ & ambulant & stationär & Zahn-Beh. & Zahn-Ers. & $\sum K^{S}_x$ & $\sum u_x$ & Summe $K_x^{\mathrm{beob}}$ \\
\midrule
\endfirsthead
\caption[]{Herleitung der Kopfschäden Frauen: beobachtete Teilkopfschäden je Leistungsbereich nach S-Kosten-Umlage (EUR p.\,a., Beobachtungsjahr 2022, Tafelbereich Alter 0--100) (Fortsetzung)}\\
\toprule
$x$ & ambulant & stationär & Zahn-Beh. & Zahn-Ers. & $\sum K^{S}_x$ & $\sum u_x$ & Summe $K_x^{\mathrm{beob}}$ \\
\midrule
\endhead
\bottomrule
\endfoot
0 & 1417,6027 & 2111,1453 & 3,1573 & 0,4190 & 0,0000 & 0,0000 & 3532,3243 \\
1 & 1182,8504 & 527,7863 & 9,0315 & 0,0000 & 0,0000 & 0,0000 & 1719,6682 \\
2 & 986,9389 & 316,6718 & 20,4028 & 0,0000 & 0,0000 & 0,0000 & 1324,0135 \\
3 & 830,3044 & 211,1145 & 37,1696 & 0,0000 & 0,0000 & 0,0000 & 1078,5885 \\
4 & 712,9282 & 158,3359 & 58,8432 & 0,0000 & 0,0000 & 0,0000 & 930,1073 \\
5 & 633,2240 & 120,5389 & 83,9979 & 7,7933 & 0,0000 & 0,0000 & 845,5541 \\
6 & 587,7256 & 116,1431 & 109,8449 & 31,9404 & 0,0000 & 0,0000 & 845,6540 \\
7 & 571,7012 & 114,1927 & 132,6315 & 79,6667 & 0,0000 & 0,0000 & 898,1921 \\
8 & 580,1448 & 117,5326 & 149,5189 & 153,4752 & 0,0000 & 0,0000 & 1000,6715 \\
9 & 608,6299 & 127,5453 & 160,7702 & 251,7286 & 0,0000 & 0,0000 & 1148,6740 \\
10 & 653,7916 & 148,9222 & 169,6180 & 388,2133 & 0,0000 & 0,0000 & 1360,5451 \\
11 & 713,2847 & 190,5666 & 179,7219 & 586,8263 & 0,0000 & 0,0000 & 1670,3995 \\
12 & 785,2391 & 254,5285 & 192,5315 & 822,3059 & 0,0000 & 0,0000 & 2054,6050 \\
13 & 867,4654 & 335,3717 & 206,4316 & 997,5789 & 0,0000 & 0,0000 & 2406,8476 \\
14 & 956,9344 & 420,4470 & 218,4875 & 1012,9752 & 0,0000 & 0,0000 & 2608,8441 \\
15 & 1049,9793 & 493,1282 & 226,3741 & 861,2275 & 0,0283 & 0,0138 & 2630,7091 \\
16 & 1143,1046 & 540,8384 & 229,3408 & 628,1579 & 0,0376 & 0,0183 & 2541,4417 \\
17 & 1233,9294 & 562,1209 & 228,2024 & 409,6379 & 0,8619 & 0,4278 & 2433,8905 \\
18 & 1321,6653 & 564,7463 & 224,4995 & 257,6958 & 0,7841 & 0,3925 & 2368,6070 \\
19 & 1406,8663 & 561,6366 & 219,8621 & 175,7385 & 1,3735 & 0,6862 & 2364,1035 \\
20 & 1490,7997 & 559,3866 & 215,4842 & 159,7979 & 3,6305 & 1,7548 & 2425,4685 \\
21 & 1575,0627 & 556,6932 & 212,2907 & 146,1427 & 8,0435 & 4,1329 & 2490,1894 \\
22 & 1661,4472 & 549,8170 & 210,9232 & 135,7855 & 15,0181 & 7,1720 & 2557,9729 \\
23 & 1748,4761 & 536,3897 & 211,5808 & 130,4824 & 31,4858 & 13,0815 & 2626,9291 \\
24 & 1834,2916 & 516,4843 & 213,9214 & 128,6573 & 64,5162 & 23,9583 & 2693,3547 \\
25 & 1919,8335 & 492,0057 & 217,2820 & 128,3677 & 119,2403 & 44,3334 & 2757,4889 \\
26 & 2001,9254 & 465,7127 & 220,9331 & 128,5558 & 201,5253 & 78,2425 & 2817,1270 \\
27 & 2071,7092 & 431,2233 & 224,1841 & 128,2508 & 312,4885 & 110,9447 & 2855,3673 \\
28 & 2122,2662 & 395,4573 & 226,6923 & 129,0864 & 444,2449 & 137,7244 & 2873,5021 \\
29 & 2158,9101 & 370,5487 & 228,4421 & 130,5762 & 586,0846 & 170,7295 & 2888,4772 \\
30 & 2173,8099 & 352,4874 & 229,6083 & 133,9517 & 726,5488 & 194,3098 & 2889,8573 \\
31 & 2182,2981 & 347,7036 & 230,5161 & 138,2492 & 854,2428 & 227,7105 & 2898,7671 \\
32 & 2167,9141 & 366,3549 & 231,4626 & 143,3442 & 956,4916 & 237,8214 & 2909,0758 \\
33 & 2155,9750 & 350,5365 & 232,7146 & 148,3559 & 1026,8575 & 253,8867 & 2887,5821 \\
34 & 2128,7609 & 335,5564 & 234,4536 & 158,4504 & 1059,0370 & 248,7523 & 2857,2215 \\
35 & 2108,9003 & 331,9225 & 236,7250 & 171,4390 & 1039,7672 & 241,7082 & 2848,9868 \\
36 & 2092,0353 & 336,2349 & 239,2484 & 184,9510 & 971,2452 & 223,6990 & 2852,4695 \\
37 & 2080,0232 & 346,2584 & 241,5194 & 199,6368 & 866,7130 & 201,2536 & 2867,4378 \\
38 & 2063,9767 & 348,7531 & 243,1182 & 205,1220 & 746,4034 & 164,0130 & 2860,9699 \\
39 & 2059,7324 & 358,4813 & 243,9242 & 217,2474 & 616,1775 & 139,2546 & 2879,3853 \\
40 & 2054,6312 & 370,8385 & 244,1263 & 228,7853 & 479,4945 & 104,4367 & 2898,3813 \\
41 & 2061,8722 & 379,7952 & 244,2490 & 239,7521 & 355,5622 & 77,7675 & 2925,6685 \\
42 & 2073,8794 & 386,7679 & 244,9671 & 250,1397 & 250,9135 & 53,2323 & 2955,7540 \\
43 & 2096,6282 & 391,2380 & 246,6956 & 259,9121 & 168,6671 & 35,3346 & 2994,4739 \\
44 & 2129,7256 & 394,2892 & 249,4042 & 269,2548 & 108,9160 & 22,9330 & 3042,6738 \\
45 & 2172,6121 & 400,2022 & 252,8725 & 278,4705 & 65,7839 & 13,5757 & 3104,1572 \\
46 & 2225,3837 & 408,1585 & 256,8796 & 287,9104 & 38,3137 & 7,9051 & 3178,3321 \\
47 & 2286,5675 & 418,3561 & 261,0560 & 297,8599 & 21,6606 & 4,3485 & 3263,8395 \\
48 & 2354,4273 & 430,7386 & 265,0939 & 308,4127 & 12,5851 & 2,5789 & 3358,6725 \\
49 & 2426,0253 & 444,8955 & 268,9777 & 319,5229 & 8,4840 & 1,7519 & 3459,4215 \\
50 & 2498,7009 & 460,1190 & 272,9131 & 331,4673 & 6,9139 & 1,3834 & 3563,2004 \\
51 & 2569,7544 & 476,2824 & 277,0328 & 344,4491 & 6,1885 & 1,2428 & 3667,5187 \\
52 & 2636,5541 & 492,8406 & 281,2809 & 358,2142 & 5,8020 & 1,1271 & 3768,8898 \\
53 & 2697,1023 & 510,2539 & 285,3680 & 372,2022 & 5,5418 & 1,1103 & 3864,9264 \\
54 & 2750,2196 & 530,1811 & 288,6346 & 385,8870 & 3,9462 & 0,7681 & 3954,9223 \\
55 & 2796,4230 & 550,9863 & 290,5027 & 398,5355 & 3,2212 & 0,6293 & 4036,4475 \\
56 & 2837,5128 & 574,3403 & 291,1638 & 410,3809 & 2,1967 & 0,4404 & 4113,3978 \\
57 & 2876,5862 & 600,9779 & 291,1858 & 422,8510 & 1,2777 & 0,2594 & 4191,6009 \\
58 & 2917,8585 & 630,8822 & 291,1622 & 437,0272 & 0,6201 & 0,1311 & 4276,9301 \\
59 & 2965,6032 & 663,6435 & 291,6228 & 453,0487 & 0,2158 & 0,0477 & 4373,9182 \\
60 & 3023,1952 & 699,5137 & 292,9953 & 470,5894 & 0,0052 & 0,0012 & 4486,2936 \\
61 & 3092,4390 & 739,9247 & 295,5222 & 488,6324 & 0,0000 & 0,0000 & 4616,5183 \\
62 & 3173,7638 & 786,1297 & 299,0635 & 506,8266 & 0,0000 & 0,0000 & 4765,7836 \\
63 & 3266,2737 & 838,4224 & 303,2019 & 524,5156 & 0,0000 & 0,0000 & 4932,4136 \\
64 & 3368,4783 & 895,9634 & 307,4659 & 540,7837 & 0,0000 & 0,0000 & 5112,6913 \\
65 & 3478,9034 & 957,8772 & 311,3240 & 555,3464 & 0,0000 & 0,0000 & 5303,4510 \\
66 & 3596,5053 & 1024,7335 & 314,2597 & 567,8152 & 0,0000 & 0,0000 & 5503,3137 \\
67 & 3720,6736 & 1097,3722 & 316,0708 & 578,3541 & 0,0000 & 0,0000 & 5712,4707 \\
68 & 3851,0130 & 1178,3693 & 316,7598 & 587,2248 & 0,0000 & 0,0000 & 5933,3669 \\
69 & 3987,1868 & 1268,4039 & 316,6461 & 594,7987 & 0,0000 & 0,0000 & 6167,0355 \\
70 & 4128,8256 & 1367,2201 & 316,3611 & 601,3374 & 0,0000 & 0,0000 & 6413,7442 \\
71 & 4275,3743 & 1475,3248 & 316,4343 & 607,0316 & 0,0000 & 0,0000 & 6674,1650 \\
72 & 4425,8354 & 1591,5965 & 317,0205 & 611,5469 & 0,0000 & 0,0000 & 6945,9993 \\
73 & 4578,4908 & 1714,0047 & 317,9353 & 614,6226 & 0,0000 & 0,0000 & 7225,0534 \\
74 & 4730,7937 & 1840,5238 & 318,5389 & 615,7965 & 0,0000 & 0,0000 & 7505,6529 \\
75 & 4879,6146 & 1968,3212 & 318,3691 & 615,1750 & 0,0000 & 0,0000 & 7781,4799 \\
76 & 5021,8262 & 2094,4296 & 317,5164 & 613,0606 & 0,0000 & 0,0000 & 8046,8328 \\
77 & 5154,6644 & 2218,3437 & 316,3148 & 609,6306 & 0,0000 & 0,0000 & 8298,9535 \\
78 & 5275,9093 & 2339,6438 & 314,9852 & 604,7918 & 0,0000 & 0,0000 & 8535,3301 \\
79 & 5383,9557 & 2458,4993 & 313,3387 & 597,8810 & 0,0000 & 0,0000 & 8753,6747 \\
80 & 5477,8313 & 2591,5745 & 310,6781 & 587,9700 & 0,0000 & 0,0000 & 8968,0539 \\
81 & 5557,1975 & 2713,9020 & 306,1552 & 574,8886 & 0,0000 & 0,0000 & 9152,1433 \\
82 & 5622,3630 & 2835,1312 & 299,2782 & 558,5112 & 0,0000 & 0,0000 & 9315,2836 \\
83 & 5674,3138 & 2951,8067 & 289,8656 & 538,8921 & 0,0000 & 0,0000 & 9454,8782 \\
84 & 5714,7217 & 3060,2610 & 277,8993 & 516,4103 & 0,0000 & 0,0000 & 9569,2923 \\
85 & 5745,8591 & 3157,7943 & 263,7329 & 491,8254 & 0,0000 & 0,0000 & 9659,2117 \\
86 & 5770,3592 & 3242,7249 & 246,6492 & 465,8729 & 0,0000 & 0,0000 & 9725,6062 \\
87 & 5790,8138 & 3314,5130 & 230,2299 & 438,8822 & 0,0000 & 0,0000 & 9774,4389 \\
88 & 5809,2978 & 3373,2394 & 213,9566 & 411,1888 & 0,0000 & 0,0000 & 9807,6826 \\
89 & 5827,0259 & 3419,4897 & 197,3140 & 383,0120 & 0,0000 & 0,0000 & 9826,8416 \\
90 & 5844,4756 & 3453,8351 & 180,9429 & 354,7409 & 0,0000 & 0,0000 & 9833,9945 \\
91 & 5861,6513 & 3476,4481 & 165,3595 & 326,5140 & 0,0000 & 0,0000 & 9829,9729 \\
92 & 5878,3313 & 3486,7165 & 150,8886 & 298,4496 & 0,0000 & 0,0000 & 9814,3860 \\
93 & 5894,2450 & 3486,7165 & 137,2414 & 270,5323 & 0,0000 & 0,0000 & 9788,7352 \\
94 & 5909,1702 & 3486,7165 & 124,7222 & 242,7729 & 0,0000 & 0,0000 & 9763,3818 \\
95 & 5922,9632 & 3486,7165 & 113,3198 & 215,1935 & 0,0000 & 0,0000 & 9738,1930 \\
96 & 5935,5499 & 3486,7165 & 102,8507 & 187,7362 & 0,0000 & 0,0000 & 9712,8533 \\
97 & 5946,9025 & 3486,7165 & 93,2937 & 160,3936 & 0,0000 & 0,0000 & 9687,3063 \\
98 & 5957,0158 & 3486,7165 & 84,4920 & 133,1526 & 0,0000 & 0,0000 & 9661,3769 \\
99 & 5965,8916 & 3486,7165 & 76,0731 & 105,9694 & 0,0000 & 0,0000 & 9634,6506 \\
100 & 5973,5308 & 3486,7165 & 67,6992 & 78,8010 & 0,0000 & 0,0000 & 9606,7475 \\
\end{longtable}
\end{small}


% ---------------------------------------------------------------------
\setzeTeil{Anlage 2.3.3}
\subsubsection{Herleitung der rechnungsmäßigen Profile und Grundkopfschäden}\label{sec:anlage233}\markright{Anlage~2.3.3}

Rechnungsmäßige Profile $k_x$ gemäß Abschnitt~\ref{sec:kopfschaden}
(Normierungsalter 40, Monotonisierung ab dem Monotonisierungsalter
$\tilde{x} = \MonoAlter$, Fortschreibung ab Alter 101) sowie
rechnungsmäßige Kopfschäden $K_x = k_x\cdot G$ mit
$G^{M} = \GMann$~EUR und $G^{F} = \GFrau$~EUR
(Gleichung~\eqref{eq:gks}). Endalter: 110.

\paragraph{Kopfschäden der Kinder und Jugendlichen.}
Die altersunabhängigen mittleren Kopfschäden gemäß
Abschnitt~\ref{sec:kopfschaden-kinder} betragen (Kinder
geschlechtsunabhängig, Jugendliche nach Geschlecht getrennt)
\begin{equation}
\bar{K}_{\mathrm{Ki}} = \KKinder, \qquad
\bar{K}^{M}_{\mathrm{Ju}} = \KJugendM, \qquad
\bar{K}^{F}_{\mathrm{Ju}} = \KJugendF
\qquad\text{(EUR p.\,a.)}.
\label{eq:kkinder}
\end{equation}

\paragraph{Festsetzung und Sicherheitsnachweis der Grundkopfschäden.}
Der rechnungsmäßige Grundkopfschaden wird je Beobachtungseinheit als
$G = K^{\mathrm{beob}}_{40}$ festgesetzt:
\begin{equation}
G^{M} = \GMann\;\text{EUR p.\,a.},
\qquad
G^{F} = \GFrau\;\text{EUR p.\,a.}
\label{eq:gks}
\end{equation}
Die rechnungsmäßigen Profile $k_x$ und die rechnungsmäßigen
Kopfschäden $K_x = k_x\cdot G$ enthält Anlage~\ref{sec:anlage233}.

\textbf{Nachweis der ausreichenden Sicherheit (Festlegung):} Auf eine
Trendextrapolation der Grundkopfschäden auf den Kalkulationsstichtag
wird verzichtet, da die rechnungsmäßigen Grundkopfschäden der
Stützdaten oberhalb der im eigenen Bestand beobachteten tatsächlichen
Grundkopfschäden liegen (Ermittlung gemäß
Abschnitt~\ref{sec:af-gks}; fiktive Werte):

\begin{center}\small
\begin{tabular}{@{}lccccc@{}}
\toprule
 & $G^{(1)}$ (2021) & $G^{(2)}$ (2022) & $G^{(3)}$ (2023) & $G^{\mathrm{ext}}$ & $G^{\mathrm{rech}}$ \\
\midrule
Männer & 1\,750{,}11 & 1\,769{,}84 & 1\,790{,}06 & 1\,829{,}93 & \GMann \\
Frauen & 2\,681{,}52 & 2\,708{,}77 & 2\,738{,}90 & 2\,795{,}80 & \GFrau \\
\bottomrule
\end{tabular}
\end{center}

Der extrapolierte tatsächliche Grundkopfschaden $G^{\mathrm{ext}}$
(Gleichung~\eqref{eq:af-ext}) unterschreitet den rechnungsmäßigen
Ansatz um 3{,}7\,\% (Männer) bzw. 3{,}5\,\% (Frauen); die
Kopfschäden sind damit ausreichend sicher im Sinne des \S\,2 Abs.\,3
KVAV. Gegenüber der vorherigen Kalkulation (Änderungsnr.~16) erhöhen
sich die rechnungsmäßigen Kopfschäden bestandsgewichtet um 9{,}0\,\%
(Männer) bzw. 7{,}0\,\% (Frauen), entsprechend den auslösenden
Faktoren der Änderungshistorie (Abschnitt~\ref{sec:historie}); die
Angemessenheit wird zudem jährlich über die auslösenden Faktoren gemäß
Abschnitt~\ref{sec:af} überprüft.


% Automatisch erzeugt von kalkulation/pipeline.py
\begin{small}
\begin{longtable}{rrr}
\caption{Rechnungsmäßige normierte Kopfschadenprofile $k_x$ (Erwachsene)}\label{tab:kx}\\
\toprule
$x$ & $k_x$ Männer & $k_x$ Frauen \\
\midrule
\endfirsthead
\caption[]{Rechnungsmäßige normierte Kopfschadenprofile $k_x$ (Erwachsene) (Fortsetzung)}\\
\toprule
$x$ & $k_x$ Männer & $k_x$ Frauen \\
\midrule
\endhead
\bottomrule
\endfoot
21 & 1,0375 & 0,8592 \\
22 & 1,0245 & 0,8826 \\
23 & 1,0064 & 0,9063 \\
24 & 0,9829 & 0,9293 \\
25 & 0,9540 & 0,9514 \\
26 & 0,9270 & 0,9720 \\
27 & 0,9411 & 0,9852 \\
28 & 0,9500 & 0,9914 \\
29 & 0,9621 & 0,9966 \\
30 & 0,9698 & 0,9971 \\
31 & 0,9840 & 1,0001 \\
32 & 0,9880 & 1,0037 \\
33 & 0,9970 & 0,9963 \\
34 & 0,9975 & 0,9858 \\
35 & 1,0010 & 0,9830 \\
36 & 1,0024 & 0,9842 \\
37 & 1,0044 & 0,9893 \\
38 & 0,9931 & 0,9871 \\
39 & 0,9976 & 0,9934 \\
40 & 1,0000 & 1,0000 \\
41 & 1,0100 & 1,0094 \\
42 & 1,0254 & 1,0198 \\
43 & 1,0486 & 1,0332 \\
44 & 1,0791 & 1,0498 \\
45 & 1,1153 & 1,0710 \\
46 & 1,1566 & 1,0966 \\
47 & 1,2017 & 1,1261 \\
48 & 1,2508 & 1,1588 \\
49 & 1,3038 & 1,1936 \\
50 & 1,3610 & 1,2294 \\
51 & 1,4221 & 1,2654 \\
52 & 1,4878 & 1,3003 \\
53 & 1,5592 & 1,3335 \\
54 & 1,6359 & 1,3645 \\
55 & 1,7179 & 1,3927 \\
56 & 1,8061 & 1,4192 \\
57 & 1,9010 & 1,4462 \\
58 & 2,0032 & 1,4756 \\
59 & 2,1130 & 1,5091 \\
60 & 2,2317 & 1,5479 \\
61 & 2,3598 & 1,5928 \\
62 & 2,4976 & 1,6443 \\
63 & 2,6448 & 1,7018 \\
64 & 2,7999 & 1,7640 \\
65 & 2,9611 & 1,8298 \\
66 & 3,1259 & 1,8988 \\
67 & 3,2918 & 1,9709 \\
68 & 3,4592 & 2,0471 \\
69 & 3,6294 & 2,1278 \\
70 & 3,8039 & 2,2129 \\
71 & 3,9822 & 2,3027 \\
72 & 4,1629 & 2,3965 \\
73 & 4,3442 & 2,4928 \\
74 & 4,5242 & 2,5896 \\
75 & 4,7004 & 2,6848 \\
76 & 4,8842 & 2,7763 \\
77 & 5,0539 & 2,8633 \\
78 & 5,2160 & 2,9449 \\
79 & 5,3684 & 3,0202 \\
80 & 5,5088 & 3,0942 \\
81 & 5,6429 & 3,1577 \\
82 & 5,7659 & 3,2140 \\
83 & 5,8734 & 3,2621 \\
84 & 5,9623 & 3,3016 \\
85 & 6,0301 & 3,3326 \\
86 & 6,0761 & 3,3555 \\
87 & 6,1019 & 3,3724 \\
88 & 6,1092 & 3,3838 \\
89 & 6,1092 & 3,3905 \\
90 & 6,1092 & 3,3929 \\
91 & 6,1092 & 3,3929 \\
92 & 6,1092 & 3,3929 \\
93 & 6,1092 & 3,3929 \\
94 & 6,1092 & 3,3929 \\
95 & 6,1092 & 3,3929 \\
96 & 6,1092 & 3,3929 \\
97 & 6,1092 & 3,3929 \\
98 & 6,1092 & 3,3929 \\
99 & 6,1092 & 3,3929 \\
100 & 6,1092 & 3,3929 \\
101 & 6,1092 & 3,3929 \\
102 & 6,1092 & 3,3929 \\
103 & 6,1092 & 3,3929 \\
104 & 6,1092 & 3,3929 \\
105 & 6,1092 & 3,3929 \\
106 & 6,1092 & 3,3929 \\
107 & 6,1092 & 3,3929 \\
108 & 6,1092 & 3,3929 \\
109 & 6,1092 & 3,3929 \\
110 & 6,1092 & 3,3929 \\
\end{longtable}
\end{small}


\begin{figure}[h!]
    \centering
    \begin{tikzpicture}
        \begin{axis}[
            width=0.92\textwidth, height=0.5\textwidth,
            grid=both, xlabel={Alter $x$},
            ylabel={normierter Kopfschaden $k_x$},
            legend pos=north west]
            % Automatisch erzeugt von kalkulation/pipeline.py
\addplot+[ffbkblau,mark=*,mark size=1pt,mark options={fill=ffbkblau}] coordinates {
(21,1.0375) (22,1.0245) (23,1.0064) (24,0.9829) (25,0.9540) (26,0.9270) (27,0.9411) (28,0.9500) (29,0.9621) (30,0.9698) (31,0.9840) (32,0.9880) (33,0.9970) (34,0.9975) (35,1.0010) (36,1.0024) (37,1.0044) (38,0.9931) (39,0.9976) (40,1.0000) (41,1.0100) (42,1.0254) (43,1.0486) (44,1.0791) (45,1.1153) (46,1.1566) (47,1.2017) (48,1.2508) (49,1.3038) (50,1.3610) (51,1.4221) (52,1.4878) (53,1.5592) (54,1.6359) (55,1.7179) (56,1.8061) (57,1.9010) (58,2.0032) (59,2.1130) (60,2.2317) (61,2.3598) (62,2.4976) (63,2.6448) (64,2.7999) (65,2.9611) (66,3.1259) (67,3.2918) (68,3.4592) (69,3.6294) (70,3.8039) (71,3.9822) (72,4.1629) (73,4.3442) (74,4.5242) (75,4.7004) (76,4.8842) (77,5.0539) (78,5.2160) (79,5.3684) (80,5.5088) (81,5.6429) (82,5.7659) (83,5.8734) (84,5.9623) (85,6.0301) (86,6.0761) (87,6.1019) (88,6.1092) (89,6.1092) (90,6.1092) (91,6.1092) (92,6.1092) (93,6.1092) (94,6.1092) (95,6.1092) (96,6.1092) (97,6.1092) (98,6.1092) (99,6.1092) (100,6.1092) (101,6.1092) (102,6.1092) (103,6.1092) (104,6.1092) (105,6.1092) (106,6.1092) (107,6.1092) (108,6.1092) (109,6.1092) (110,6.1092)};
\addlegendentry{M\"anner}
\addplot+[ffbkrot,mark=*,mark size=1pt,mark options={fill=ffbkrot}] coordinates {
(21,0.8592) (22,0.8826) (23,0.9063) (24,0.9293) (25,0.9514) (26,0.9720) (27,0.9852) (28,0.9914) (29,0.9966) (30,0.9971) (31,1.0001) (32,1.0037) (33,0.9963) (34,0.9858) (35,0.9830) (36,0.9842) (37,0.9893) (38,0.9871) (39,0.9934) (40,1.0000) (41,1.0094) (42,1.0198) (43,1.0332) (44,1.0498) (45,1.0710) (46,1.0966) (47,1.1261) (48,1.1588) (49,1.1936) (50,1.2294) (51,1.2654) (52,1.3003) (53,1.3335) (54,1.3645) (55,1.3927) (56,1.4192) (57,1.4462) (58,1.4756) (59,1.5091) (60,1.5479) (61,1.5928) (62,1.6443) (63,1.7018) (64,1.7640) (65,1.8298) (66,1.8988) (67,1.9709) (68,2.0471) (69,2.1278) (70,2.2129) (71,2.3027) (72,2.3965) (73,2.4928) (74,2.5896) (75,2.6848) (76,2.7763) (77,2.8633) (78,2.9449) (79,3.0202) (80,3.0942) (81,3.1577) (82,3.2140) (83,3.2621) (84,3.3016) (85,3.3326) (86,3.3555) (87,3.3724) (88,3.3838) (89,3.3905) (90,3.3929) (91,3.3929) (92,3.3929) (93,3.3929) (94,3.3929) (95,3.3929) (96,3.3929) (97,3.3929) (98,3.3929) (99,3.3929) (100,3.3929) (101,3.3929) (102,3.3929) (103,3.3929) (104,3.3929) (105,3.3929) (106,3.3929) (107,3.3929) (108,3.3929) (109,3.3929) (110,3.3929)};
\addlegendentry{Frauen}

        \end{axis}
    \end{tikzpicture}
    \caption{Rechnungsmäßige Profile des Tarifs VitalPlus}
    \label{fig:profile}
\end{figure}

\clearpage
% Automatisch erzeugt von kalkulation/pipeline.py
\begin{small}
\begin{longtable}{rrr}
\caption{Rechnungsmäßige Kopfschäden $K_x = k_x \cdot G$ (EUR p.\,a.)}\label{tab:Kx}\\
\toprule
$x$ & $K_x$ Männer & $K_x$ Frauen \\
\midrule
\endfirsthead
\caption[]{Rechnungsmäßige Kopfschäden $K_x = k_x \cdot G$ (EUR p.\,a.) (Fortsetzung)}\\
\toprule
$x$ & $K_x$ Männer & $K_x$ Frauen \\
\midrule
\endhead
\bottomrule
\endfoot
21 & 1972,2948 & 2490,2892 \\
22 & 1947,5817 & 2558,1113 \\
23 & 1913,1734 & 2626,8030 \\
24 & 1868,4998 & 2693,4657 \\
25 & 1813,5607 & 2757,5200 \\
26 & 1762,2335 & 2817,2266 \\
27 & 1789,0377 & 2855,4853 \\
28 & 1805,9567 & 2873,4552 \\
29 & 1828,9588 & 2888,5268 \\
30 & 1843,5966 & 2889,9760 \\
31 & 1870,5909 & 2898,6711 \\
32 & 1878,1949 & 2909,1053 \\
33 & 1895,3040 & 2887,6573 \\
34 & 1896,2545 & 2857,2243 \\
35 & 1902,9080 & 2849,1088 \\
36 & 1905,5694 & 2852,5869 \\
37 & 1909,3714 & 2867,3686 \\
38 & 1887,8901 & 2860,9922 \\
39 & 1896,4446 & 2879,2520 \\
40 & 1901,0070 & 2898,3813 \\
41 & 1920,0171 & 2925,6261 \\
42 & 1949,2926 & 2955,7692 \\
43 & 1993,3959 & 2994,6076 \\
44 & 2051,3767 & 3042,7207 \\
45 & 2120,1931 & 3104,1664 \\
46 & 2198,7047 & 3178,3649 \\
47 & 2284,4401 & 3263,8672 \\
48 & 2377,7796 & 3358,6443 \\
49 & 2478,5329 & 3459,5079 \\
50 & 2587,2705 & 3563,2700 \\
51 & 2703,4221 & 3667,6117 \\
52 & 2828,3182 & 3768,7652 \\
53 & 2964,0501 & 3864,9915 \\
54 & 3109,8574 & 3954,8413 \\
55 & 3265,7399 & 4036,5756 \\
56 & 3433,4087 & 4113,3827 \\
57 & 3613,8143 & 4191,6390 \\
58 & 3808,0972 & 4276,8514 \\
59 & 4016,8278 & 4373,9472 \\
60 & 4242,4773 & 4486,4044 \\
61 & 4485,9963 & 4616,5417 \\
62 & 4747,9551 & 4765,8084 \\
63 & 5027,7833 & 4932,4653 \\
64 & 5322,6295 & 5112,7446 \\
65 & 5629,0718 & 5303,4581 \\
66 & 5942,3578 & 5503,4464 \\
67 & 6257,7348 & 5712,4197 \\
68 & 6575,9634 & 5933,2764 \\
69 & 6899,5148 & 6167,1757 \\
70 & 7231,2405 & 6413,8280 \\
71 & 7570,1901 & 6674,1026 \\
72 & 7913,7020 & 6945,9708 \\
73 & 8258,3546 & 7225,0849 \\
74 & 8600,5359 & 7505,6482 \\
75 & 8935,4933 & 7781,5741 \\
76 & 9284,8984 & 8046,7760 \\
77 & 9607,4993 & 8298,9352 \\
78 & 9915,6525 & 8535,4431 \\
79 & 10205,3660 & 8753,6912 \\
80 & 10472,2674 & 8968,1714 \\
81 & 10727,1924 & 9152,2186 \\
82 & 10961,0163 & 9315,3975 \\
83 & 11165,3745 & 9454,8096 \\
84 & 11334,3740 & 9569,2957 \\
85 & 11463,2623 & 9659,1455 \\
86 & 11550,7086 & 9725,5185 \\
87 & 11599,7546 & 9774,5011 \\
88 & 11613,6320 & 9807,5426 \\
89 & 11613,6320 & 9826,9618 \\
90 & 11613,6320 & 9833,9179 \\
91 & 11613,6320 & 9833,9179 \\
92 & 11613,6320 & 9833,9179 \\
93 & 11613,6320 & 9833,9179 \\
94 & 11613,6320 & 9833,9179 \\
95 & 11613,6320 & 9833,9179 \\
96 & 11613,6320 & 9833,9179 \\
97 & 11613,6320 & 9833,9179 \\
98 & 11613,6320 & 9833,9179 \\
99 & 11613,6320 & 9833,9179 \\
100 & 11613,6320 & 9833,9179 \\
101 & 11613,6320 & 9833,9179 \\
102 & 11613,6320 & 9833,9179 \\
103 & 11613,6320 & 9833,9179 \\
104 & 11613,6320 & 9833,9179 \\
105 & 11613,6320 & 9833,9179 \\
106 & 11613,6320 & 9833,9179 \\
107 & 11613,6320 & 9833,9179 \\
108 & 11613,6320 & 9833,9179 \\
109 & 11613,6320 & 9833,9179 \\
110 & 11613,6320 & 9833,9179 \\
\end{longtable}
\end{small}


\begin{figure}[h!]
    \centering
    \begin{tikzpicture}
        \begin{axis}[
            width=0.92\textwidth, height=0.5\textwidth,
            grid=both, xlabel={Alter $x$},
            ylabel={rechnungsmäßiger Kopfschaden $K_x$ (EUR p.\,a.)},
            legend pos=north west]
            % Automatisch erzeugt von kalkulation/pipeline.py
\addplot+[ffbkblau,mark=*,mark size=1pt,mark options={fill=ffbkblau}] coordinates {
(21,1972.2948) (22,1947.5817) (23,1913.1734) (24,1868.4998) (25,1813.5607) (26,1762.2335) (27,1789.0377) (28,1805.9567) (29,1828.9588) (30,1843.5966) (31,1870.5909) (32,1878.1949) (33,1895.3040) (34,1896.2545) (35,1902.9080) (36,1905.5694) (37,1909.3714) (38,1887.8901) (39,1896.4446) (40,1901.0070) (41,1920.0171) (42,1949.2926) (43,1993.3959) (44,2051.3767) (45,2120.1931) (46,2198.7047) (47,2284.4401) (48,2377.7796) (49,2478.5329) (50,2587.2705) (51,2703.4221) (52,2828.3182) (53,2964.0501) (54,3109.8574) (55,3265.7399) (56,3433.4087) (57,3613.8143) (58,3808.0972) (59,4016.8278) (60,4242.4773) (61,4485.9963) (62,4747.9551) (63,5027.7833) (64,5322.6295) (65,5629.0718) (66,5942.3578) (67,6257.7348) (68,6575.9634) (69,6899.5148) (70,7231.2405) (71,7570.1901) (72,7913.7020) (73,8258.3546) (74,8600.5359) (75,8935.4933) (76,9284.8984) (77,9607.4993) (78,9915.6525) (79,10205.3660) (80,10472.2674) (81,10727.1924) (82,10961.0163) (83,11165.3745) (84,11334.3740) (85,11463.2623) (86,11550.7086) (87,11599.7546) (88,11613.6320) (89,11613.6320) (90,11613.6320) (91,11613.6320) (92,11613.6320) (93,11613.6320) (94,11613.6320) (95,11613.6320) (96,11613.6320) (97,11613.6320) (98,11613.6320) (99,11613.6320) (100,11613.6320) (101,11613.6320) (102,11613.6320) (103,11613.6320) (104,11613.6320) (105,11613.6320) (106,11613.6320) (107,11613.6320) (108,11613.6320) (109,11613.6320) (110,11613.6320)};
\addlegendentry{M\"anner}
\addplot+[ffbkrot,mark=*,mark size=1pt,mark options={fill=ffbkrot}] coordinates {
(21,2490.2892) (22,2558.1113) (23,2626.8030) (24,2693.4657) (25,2757.5200) (26,2817.2266) (27,2855.4853) (28,2873.4552) (29,2888.5268) (30,2889.9760) (31,2898.6711) (32,2909.1053) (33,2887.6573) (34,2857.2243) (35,2849.1088) (36,2852.5869) (37,2867.3686) (38,2860.9922) (39,2879.2520) (40,2898.3813) (41,2925.6261) (42,2955.7692) (43,2994.6076) (44,3042.7207) (45,3104.1664) (46,3178.3649) (47,3263.8672) (48,3358.6443) (49,3459.5079) (50,3563.2700) (51,3667.6117) (52,3768.7652) (53,3864.9915) (54,3954.8413) (55,4036.5756) (56,4113.3827) (57,4191.6390) (58,4276.8514) (59,4373.9472) (60,4486.4044) (61,4616.5417) (62,4765.8084) (63,4932.4653) (64,5112.7446) (65,5303.4581) (66,5503.4464) (67,5712.4197) (68,5933.2764) (69,6167.1757) (70,6413.8280) (71,6674.1026) (72,6945.9708) (73,7225.0849) (74,7505.6482) (75,7781.5741) (76,8046.7760) (77,8298.9352) (78,8535.4431) (79,8753.6912) (80,8968.1714) (81,9152.2186) (82,9315.3975) (83,9454.8096) (84,9569.2957) (85,9659.1455) (86,9725.5185) (87,9774.5011) (88,9807.5426) (89,9826.9618) (90,9833.9179) (91,9833.9179) (92,9833.9179) (93,9833.9179) (94,9833.9179) (95,9833.9179) (96,9833.9179) (97,9833.9179) (98,9833.9179) (99,9833.9179) (100,9833.9179) (101,9833.9179) (102,9833.9179) (103,9833.9179) (104,9833.9179) (105,9833.9179) (106,9833.9179) (107,9833.9179) (108,9833.9179) (109,9833.9179) (110,9833.9179)};
\addlegendentry{Frauen}

        \end{axis}
    \end{tikzpicture}
    \caption{Rechnungsmäßige Kopfschäden $K_x = k_x\cdot G$ des Tarifs
    VitalPlus (Männer und Frauen)}
    \label{fig:kopfschaeden}
\end{figure}

% ---------------------------------------------------------------------
\setzeTeil{Anlage 2.4}
\subsection{Zuschläge}\label{sec:anlage241}
\markright{Anlage~2.4}
Der Sicherheitszuschlag ist festgesetzt auf
\begin{equation}
\sigma = \SigmaZ\,\%
\label{eq:sigma}
\end{equation}
des Bruttobeitrags für alle Beobachtungseinheiten und Alter.


\begin{center}\small
\begin{tabular}{@{}l*{4}{>{\centering\arraybackslash}p{2.1cm}}@{}}
\toprule
 & \multicolumn{2}{c}{Männer} & \multicolumn{2}{c}{Frauen} \\
\cmidrule(lr){2-3}\cmidrule(lr){4-5}
Zuschlag & $x\le 64$ & $x\ge 65$ & $x\le 64$ & $x\ge 65$ \\
\midrule
Sicherheitszuschlag $\sigma$ & \SigmaZ\,\% & \SigmaZ\,\% & \SigmaZ\,\% & \SigmaZ\,\% \\
Standard-/Basistarif $\Omega_x$ & \OmegaU\,\% & \OmegaO\,\% & \OmegaU\,\% & \OmegaO\,\% \\
Kostensätze $\Delta_x$ (rechnerisch) & \DeltaU\,\% & \DeltaO\,\% & \DeltaU\,\% & \DeltaO\,\% \\
\midrule
$\gamma^{\mathrm{Kosten}}_x$ (EUR p.\,a.) & \GammaKostMu & \GammaKostMo & \GammaKostFu & \GammaKostFo \\
$\gamma^{\mathrm{euBR}}$ (EUR p.\,a.) & \GammaEuBRM & \GammaEuBRM & \GammaEuBRF & \GammaEuBRF \\
$\gamma_x$ gesamt (EUR p.\,a.) & \GammaMu & \GammaMo & \GammaFu & \GammaFo \\
\midrule
Zillmersatz $\alpha_x$ (Monatsbeiträge) & \multicolumn{4}{c}{$\AlphaZ$ für $21\le x\le 75$, abschmelzend bis 0 ab Alter 80 (Gl.~\eqref{eq:zillmer})} \\
\bottomrule
\end{tabular}
\end{center}

\textbf{Kostensätze (Rahmen-TB \mbox{RTB-Kosten}, Version 2024-2)} in
Prozent des Bruttobeitrags; sie sind Grundlage der rechnerischen Summen
$\Delta_x$ der obigen Übersicht und der Umrechnung in die absoluten
Kostenzuschläge gemäß Abschnitt~\ref{sec:kosten}:

\begin{center}\small
\begin{tabular}{@{}lcc@{}}
\toprule
Kostenart (\S\,8 Abs.\,1 KVAV) & $x \le 64$ & $x \ge 65$ \\
\midrule
unmittelbare Abschlusskosten, laufender Zuschlag $\alpha_u$ & 2{,}0\,\% & --- \\
mittelbare Abschlusskosten $\alpha_m$ & 5{,}0\,\% & 5{,}0\,\% \\
Schadenregulierungskosten $\rho$ & 4{,}4\,\% & 4{,}4\,\% \\
sonstige Verwaltungskosten $\beta$ & 2{,}6\,\% & 2{,}6\,\% \\
\midrule
Summe $\Delta_x$ & \DeltaU\,\% & \DeltaO\,\% \\
\bottomrule
\end{tabular}
\end{center}

Kinder/Jugendliche: $\gamma^{K} = \GammaKi$~EUR p.\,a.
($\gamma^{\mathrm{euBR}} = \GammaEuBRK$~EUR,
$\gamma^{\mathrm{Kosten}} = \GammaKostKi$~EUR); Gamma-Alter:
Männer \GammaAlterM, Frauen \GammaAlterF. Die Basistarif-Komponente
von $\Omega_x$ (0{,}6\,\%) ist branchenweit einheitlich, die
Standardtarif-Komponente (0{,}1\,\%, $x \le 64$) poolbedingt
(Abschnitt~\ref{sec:omega}).

% ---------------------------------------------------------------------
\setzeTeil{Anlage 3}
\setcounter{section}{2}
\section{Berechnung der Beiträge für den Neuzugang}\label{sec:anlagenteil3}

Anlagen zu Abschnitt~3: versicherungstechnische Grundwerte.

\setzeTeil{Anlage 3.1}
\subsection{Versicherungstechnische Grundwerte}\label{sec:anlage311}

Leistungsbarwerte $A_x$ (Profileinheiten) und Rentenbarwerte $a_x$
gemäß den Rekursionen~\eqref{eq:rekA} und \eqref{eq:reka}
(Rechnungszins: Männer \ZinsM\,\%, Frauen \ZinsF\,\%). Endalter: 110.

% Automatisch erzeugt von kalkulation/pipeline.py
\begin{small}
\begin{longtable}{rrrrr}
\caption{Leistungsbarwerte $A_x$ und Rentenbarwerte $a_x$ (in Profileinheiten)}\label{tab:barwerte}\\
\toprule
$x$ & $A_x$ M & $a_x$ M & $A_x$ F & $a_x$ F \\
\midrule
\endfirsthead
\caption[]{Leistungsbarwerte $A_x$ und Rentenbarwerte $a_x$ (in Profileinheiten) (Fortsetzung)}\\
\toprule
$x$ & $A_x$ M & $a_x$ M & $A_x$ F & $a_x$ F \\
\midrule
\endhead
\bottomrule
\endfoot
21 & 12,6404 & 9,8230 & 12,6073 & 10,8979 \\
22 & 13,3198 & 10,1286 & 13,3919 & 11,2828 \\
23 & 14,2746 & 10,5981 & 14,4208 & 11,8541 \\
24 & 15,5805 & 11,2708 & 15,7576 & 12,6557 \\
25 & 16,9858 & 11,9511 & 17,1328 & 13,4671 \\
26 & 18,4868 & 12,6280 & 18,5282 & 14,2752 \\
27 & 20,0680 & 13,2889 & 19,9235 & 15,0653 \\
28 & 21,6654 & 13,9199 & 21,3023 & 15,8210 \\
29 & 23,2590 & 14,5063 & 22,6462 & 16,5251 \\
30 & 24,8441 & 15,0493 & 23,9554 & 17,1786 \\
31 & 26,4010 & 15,5362 & 25,2120 & 17,7667 \\
32 & 27,8966 & 15,9543 & 26,3897 & 18,2748 \\
33 & 29,3146 & 16,2914 & 27,4642 & 18,6890 \\
34 & 30,6223 & 16,5359 & 28,4238 & 18,9961 \\
35 & 31,9688 & 16,7652 & 29,4038 & 19,2855 \\
36 & 33,3486 & 16,9772 & 30,3935 & 19,5547 \\
37 & 34,7605 & 17,1697 & 31,3853 & 19,8014 \\
38 & 36,2005 & 17,3406 & 32,3712 & 20,0232 \\
39 & 37,6792 & 17,4878 & 33,3547 & 20,2177 \\
40 & 39,2173 & 17,6276 & 34,3585 & 20,4037 \\
41 & 40,8183 & 17,7592 & 35,3815 & 20,5804 \\
42 & 42,4753 & 17,8820 & 36,4201 & 20,7470 \\
43 & 44,1833 & 17,9953 & 37,4723 & 20,9028 \\
44 & 45,9342 & 18,0983 & 38,5337 & 21,0468 \\
45 & 47,7684 & 18,2088 & 39,6403 & 21,2000 \\
46 & 49,6859 & 18,3275 & 40,7900 & 21,3631 \\
47 & 51,6878 & 18,4552 & 41,9812 & 21,5368 \\
48 & 53,7772 & 18,5931 & 43,2127 & 21,7219 \\
49 & 55,9573 & 18,7422 & 44,4839 & 21,9193 \\
50 & 58,2318 & 18,9038 & 45,7955 & 22,1299 \\
51 & 60,6045 & 19,0792 & 47,1495 & 22,3547 \\
52 & 63,0795 & 19,2697 & 48,5487 & 22,5948 \\
53 & 65,6601 & 19,4765 & 49,9974 & 22,8512 \\
54 & 68,3496 & 19,7012 & 51,5011 & 23,1252 \\
55 & 70,6831 & 19,8139 & 52,7162 & 23,2636 \\
56 & 72,6089 & 19,8079 & 53,6180 & 23,2589 \\
57 & 74,0808 & 19,6787 & 54,1859 & 23,1062 \\
58 & 75,0589 & 19,4238 & 54,4029 & 22,8033 \\
59 & 75,5110 & 19,0430 & 54,2564 & 22,3508 \\
60 & 75,8844 & 18,6542 & 54,0738 & 21,8877 \\
61 & 76,1618 & 18,2556 & 53,8504 & 21,4144 \\
62 & 76,3302 & 17,8467 & 53,5806 & 20,9312 \\
63 & 76,3777 & 17,4274 & 53,2583 & 20,4385 \\
64 & 76,2894 & 16,9970 & 52,8775 & 19,9366 \\
65 & 76,0576 & 16,5560 & 52,4291 & 19,4240 \\
66 & 75,6774 & 16,1053 & 51,9067 & 18,9001 \\
67 & 75,1444 & 15,6451 & 51,3059 & 18,3647 \\
68 & 74,4553 & 15,1756 & 50,6229 & 17,8180 \\
69 & 73,6072 & 14,6969 & 49,8539 & 17,2605 \\
70 & 72,6076 & 14,2116 & 49,0049 & 16,6962 \\
71 & 71,4530 & 13,7203 & 48,0706 & 16,1251 \\
72 & 70,1417 & 13,2238 & 47,0460 & 15,5475 \\
73 & 68,6723 & 12,7228 & 45,9272 & 14,9638 \\
74 & 67,0488 & 12,2186 & 44,7133 & 14,3750 \\
75 & 65,2783 & 11,7127 & 43,4069 & 13,7825 \\
76 & 63,3700 & 11,2065 & 42,0134 & 13,1878 \\
77 & 61,3142 & 10,7001 & 40,5411 & 12,5928 \\
78 & 59,1109 & 10,1916 & 38,9888 & 11,9962 \\
79 & 56,7749 & 9,6828 & 37,3621 & 11,3984 \\
80 & 54,3261 & 9,1759 & 35,6749 & 10,8020 \\
81 & 51,7890 & 8,6736 & 33,9372 & 10,2101 \\
82 & 49,1837 & 8,1787 & 32,1689 & 9,6258 \\
83 & 46,5347 & 7,6940 & 30,3840 & 9,0515 \\
84 & 43,8652 & 7,2215 & 28,6012 & 8,4906 \\
85 & 41,2092 & 6,7642 & 26,8392 & 7,9464 \\
86 & 38,6026 & 6,3252 & 25,1163 & 7,4221 \\
87 & 36,0812 & 5,9072 & 23,4495 & 6,9205 \\
88 & 33,6784 & 5,5127 & 21,8528 & 6,4441 \\
89 & 31,4256 & 5,1439 & 20,3387 & 5,9952 \\
90 & 29,3424 & 4,8029 & 18,9184 & 5,5759 \\
91 & 27,4405 & 4,4916 & 17,6021 & 5,1880 \\
92 & 25,7323 & 4,2120 & 16,3982 & 4,8332 \\
93 & 24,2328 & 3,9666 & 15,3129 & 4,5133 \\
94 & 22,9536 & 3,7572 & 14,3514 & 4,2299 \\
95 & 21,9219 & 3,5883 & 13,4782 & 3,9725 \\
96 & 21,1738 & 3,4659 & 12,6927 & 3,7410 \\
97 & 20,4715 & 3,3510 & 11,9953 & 3,5354 \\
98 & 19,7270 & 3,2291 & 11,3882 & 3,3565 \\
99 & 18,8983 & 3,0934 & 10,8770 & 3,2058 \\
100 & 17,9238 & 2,9339 & 10,3082 & 3,0382 \\
101 & 16,7134 & 2,7358 & 9,5834 & 2,8246 \\
102 & 15,1328 & 2,4771 & 8,6392 & 2,5463 \\
103 & 12,9818 & 2,1250 & 7,3696 & 2,1721 \\
104 & 11,3021 & 1,8500 & 6,3855 & 1,8820 \\
105 & 9,9657 & 1,6313 & 5,6077 & 1,6528 \\
106 & 8,8845 & 1,4543 & 4,9821 & 1,4684 \\
107 & 7,9965 & 1,3089 & 4,4708 & 1,3177 \\
108 & 7,2573 & 1,1879 & 4,0472 & 1,1928 \\
109 & 6,6348 & 1,0860 & 3,6919 & 1,0881 \\
110 & 6,1092 & 1,0000 & 3,3929 & 1,0000 \\
\end{longtable}
\end{small}


% ---------------------------------------------------------------------
\setzeTeil{Anlage 5}
\setcounter{section}{4}
\section{Ermittlung der Alterungsrückstellung}\label{sec:anlagenteil5}

Anlagen zu Abschnitt~5: gemittelte Barwerte für die
Bilanzdeckungsrückstellung und Nachweis nach \S\,8 Abs.\,3 KVAV.

\setzeTeil{Anlage 5.2}
\setcounter{subsection}{1}
\subsection{Bilanzdeckungsrückstellung}\label{sec:anlage521}

\phantomsection\label{sec:anlage522}%
Gemittelte Barwerte gemäß Gleichung~\eqref{eq:halbbarwerte} für die
Bilanzdeckungsrückstellung (Abschnitt~\ref{sec:bilanzar}) sowie der
Nachweis nach \S\,8 Abs.\,3 KVAV. Endalter: 110.

% Automatisch erzeugt von kalkulation/pipeline.py
\begin{small}
\begin{longtable}{rrrrr}
\caption{Gemittelte Barwerte $A_{x+1/2}$, $a_{x+1/2}$ für die Bilanzdeckungsrückstellung}\label{tab:barwertehalb}\\
\toprule
$x$ & $A_{x+1/2}$ M & $a_{x+1/2}$ M & $A_{x+1/2}$ F & $a_{x+1/2}$ F \\
\midrule
\endfirsthead
\caption[]{Gemittelte Barwerte $A_{x+1/2}$, $a_{x+1/2}$ für die Bilanzdeckungsrückstellung (Fortsetzung)}\\
\toprule
$x$ & $A_{x+1/2}$ M & $a_{x+1/2}$ M & $A_{x+1/2}$ F & $a_{x+1/2}$ F \\
\midrule
\endhead
\bottomrule
\endfoot
21 & 12,9801 & 9,9758 & 12,9996 & 11,0904 \\
22 & 13,7972 & 10,3634 & 13,9064 & 11,5685 \\
23 & 14,9276 & 10,9345 & 15,0892 & 12,2549 \\
24 & 16,2832 & 11,6109 & 16,4452 & 13,0614 \\
25 & 17,7363 & 12,2896 & 17,8305 & 13,8712 \\
26 & 19,2774 & 12,9585 & 19,2259 & 14,6703 \\
27 & 20,8667 & 13,6044 & 20,6129 & 15,4432 \\
28 & 22,4622 & 14,2131 & 21,9742 & 16,1731 \\
29 & 24,0516 & 14,7778 & 23,3008 & 16,8519 \\
30 & 25,6226 & 15,2928 & 24,5837 & 17,4727 \\
31 & 27,1488 & 15,7452 & 25,8009 & 18,0208 \\
32 & 28,6056 & 16,1229 & 26,9270 & 18,4819 \\
33 & 29,9684 & 16,4137 & 27,9440 & 18,8426 \\
34 & 31,2956 & 16,6506 & 28,9138 & 19,1408 \\
35 & 32,6587 & 16,8712 & 29,8987 & 19,4201 \\
36 & 34,0546 & 17,0735 & 30,8894 & 19,6781 \\
37 & 35,4805 & 17,2552 & 31,8783 & 19,9123 \\
38 & 36,9399 & 17,4142 & 32,8630 & 20,1204 \\
39 & 38,4483 & 17,5577 & 33,8566 & 20,3107 \\
40 & 40,0178 & 17,6934 & 34,8700 & 20,4921 \\
41 & 41,6468 & 17,8206 & 35,9008 & 20,6637 \\
42 & 43,3293 & 17,9387 & 36,9462 & 20,8249 \\
43 & 45,0588 & 18,0468 & 38,0030 & 20,9748 \\
44 & 46,8513 & 18,1536 & 39,0870 & 21,1234 \\
45 & 48,7271 & 18,2682 & 40,2152 & 21,2816 \\
46 & 50,6869 & 18,3914 & 41,3856 & 21,4500 \\
47 & 52,7325 & 18,5242 & 42,5970 & 21,6294 \\
48 & 54,8673 & 18,6677 & 43,8483 & 21,8206 \\
49 & 57,0946 & 18,8230 & 45,1397 & 22,0246 \\
50 & 59,4182 & 18,9915 & 46,4725 & 22,2423 \\
51 & 61,8420 & 19,1745 & 47,8491 & 22,4748 \\
52 & 64,3698 & 19,3731 & 49,2731 & 22,7230 \\
53 & 67,0049 & 19,5889 & 50,7493 & 22,9882 \\
54 & 69,5163 & 19,7576 & 52,1087 & 23,1944 \\
55 & 71,6460 & 19,8109 & 53,1671 & 23,2613 \\
56 & 73,3449 & 19,7433 & 53,9020 & 23,1826 \\
57 & 74,5699 & 19,5513 & 54,2944 & 22,9548 \\
58 & 75,2850 & 19,2334 & 54,3297 & 22,5771 \\
59 & 75,6977 & 18,8486 & 54,1651 & 22,1193 \\
60 & 76,0231 & 18,4549 & 53,9621 & 21,6510 \\
61 & 76,2460 & 18,0512 & 53,7155 & 21,1728 \\
62 & 76,3540 & 17,6371 & 53,4195 & 20,6849 \\
63 & 76,3336 & 17,2122 & 53,0679 & 20,1876 \\
64 & 76,1735 & 16,7765 & 52,6533 & 19,6803 \\
65 & 75,8675 & 16,3307 & 52,1679 & 19,1621 \\
66 & 75,4109 & 15,8752 & 51,6063 & 18,6324 \\
67 & 74,7998 & 15,4104 & 50,9644 & 18,0914 \\
68 & 74,0313 & 14,9363 & 50,2384 & 17,5393 \\
69 & 73,1074 & 14,4543 & 49,4294 & 16,9784 \\
70 & 72,0303 & 13,9660 & 48,5378 & 16,4107 \\
71 & 70,7974 & 13,4721 & 47,5583 & 15,8363 \\
72 & 69,4070 & 12,9733 & 46,4866 & 15,2557 \\
73 & 67,8606 & 12,4707 & 45,3203 & 14,6694 \\
74 & 66,1636 & 11,9657 & 44,0601 & 14,0788 \\
75 & 64,3242 & 11,4596 & 42,7102 & 13,4852 \\
76 & 62,3421 & 10,9533 & 41,2772 & 12,8903 \\
77 & 60,2126 & 10,4459 & 39,7650 & 12,2945 \\
78 & 57,9429 & 9,9372 & 38,1755 & 11,6973 \\
79 & 55,5505 & 9,4294 & 36,5185 & 11,1002 \\
80 & 53,0576 & 8,9248 & 34,8061 & 10,5061 \\
81 & 50,4864 & 8,4262 & 33,0531 & 9,9180 \\
82 & 47,8592 & 7,9363 & 31,2765 & 9,3387 \\
83 & 45,2000 & 7,4578 & 29,4926 & 8,7711 \\
84 & 42,5372 & 6,9929 & 27,7202 & 8,2185 \\
85 & 39,9059 & 6,5447 & 25,9778 & 7,6843 \\
86 & 37,3419 & 6,1162 & 24,2829 & 7,1713 \\
87 & 34,8798 & 5,7099 & 22,6512 & 6,6823 \\
88 & 32,5520 & 5,3283 & 21,0958 & 6,2197 \\
89 & 30,3840 & 4,9734 & 19,6285 & 5,7856 \\
90 & 28,3915 & 4,6473 & 18,2603 & 5,3820 \\
91 & 26,5864 & 4,3518 & 17,0001 & 5,0106 \\
92 & 24,9826 & 4,0893 & 15,8556 & 4,6732 \\
93 & 23,5932 & 3,8619 & 14,8322 & 4,3716 \\
94 & 22,4378 & 3,6727 & 13,9148 & 4,1012 \\
95 & 21,5479 & 3,5271 & 13,0855 & 3,8568 \\
96 & 20,8227 & 3,4085 & 12,3440 & 3,6382 \\
97 & 20,0992 & 3,2901 & 11,6917 & 3,4460 \\
98 & 19,3126 & 3,1613 & 11,1326 & 3,2812 \\
99 & 18,4111 & 3,0137 & 10,5926 & 3,1220 \\
100 & 17,3186 & 2,8349 & 9,9458 & 2,9314 \\
101 & 15,9231 & 2,6064 & 9,1113 & 2,6855 \\
102 & 14,0573 & 2,3011 & 8,0044 & 2,3592 \\
103 & 12,1420 & 1,9875 & 6,8776 & 2,0271 \\
104 & 10,6339 & 1,7407 & 5,9966 & 1,7674 \\
105 & 9,4251 & 1,5428 & 5,2949 & 1,5606 \\
106 & 8,4405 & 1,3816 & 4,7265 & 1,3931 \\
107 & 7,6269 & 1,2484 & 4,2590 & 1,2553 \\
108 & 6,9461 & 1,1370 & 3,8696 & 1,1405 \\
109 & 6,3720 & 1,0430 & 3,5424 & 1,0441 \\
110 & 3,0546 & 0,5000 & 1,6965 & 0,5000 \\
\end{longtable}
\end{small}


\clearpage
\textbf{Nachweis nach \S\,8 Abs.\,3 KVAV:} Für jedes Eintrittsalter
$x_e$ ist das erste Versicherungsjahr $m$ angegeben, ab dem die
gezillmerte Alterungsrückstellung $V_{x_e+m}(b_{x_e})$ gemäß
Gleichung~\eqref{eq:ar} nichtnegativ ist (Neuzugangsbeitrag, ohne
Einmalbeträge; Eintrittsalter 21 bis 75, Endalter der
Rückstellungsbewertung: 110). Das Maximum beträgt \MaxNegM~Jahre
(Männer) bzw. \MaxNegF~Jahre (Frauen); die Grenzen des \S\,8 Abs.\,3
KVAV (vier Jahre je Zugangsjahrgang, fünfzehn Jahre und halbe künftige
Vertragsdauer je Einzelrückstellung) sind eingehalten.

% Automatisch erzeugt von kalkulation/pipeline.py
\begin{small}
\begin{longtable}{rrr}
\caption{Nachweis nach \S\,8 Abs.\,3 KVAV: erstes Vertragsjahr $m$ mit nichtnegativer gezillmerter Alterungsrückstellung, je Eintrittsalter}\label{tab:negativ}\\
\toprule
$x_e$ & $m$ Männer & $m$ Frauen \\
\midrule
\endfirsthead
\caption[]{Nachweis nach \S\,8 Abs.\,3 KVAV: erstes Vertragsjahr $m$ mit nichtnegativer gezillmerter Alterungsrückstellung, je Eintrittsalter (Fortsetzung)}\\
\toprule
$x_e$ & $m$ Männer & $m$ Frauen \\
\midrule
\endhead
\bottomrule
\endfoot
21 & 3 & 2 \\
22 & 2 & 2 \\
23 & 2 & 2 \\
24 & 2 & 2 \\
25 & 2 & 2 \\
26 & 2 & 2 \\
27 & 2 & 2 \\
28 & 2 & 2 \\
29 & 2 & 2 \\
30 & 1 & 2 \\
31 & 1 & 2 \\
32 & 1 & 2 \\
33 & 1 & 2 \\
34 & 1 & 2 \\
35 & 1 & 2 \\
36 & 1 & 2 \\
37 & 1 & 2 \\
38 & 1 & 1 \\
39 & 1 & 1 \\
40 & 1 & 1 \\
41 & 1 & 1 \\
42 & 1 & 1 \\
43 & 1 & 1 \\
44 & 1 & 1 \\
45 & 1 & 1 \\
46 & 1 & 1 \\
47 & 1 & 1 \\
48 & 1 & 1 \\
49 & 1 & 1 \\
50 & 1 & 1 \\
51 & 1 & 1 \\
52 & 1 & 1 \\
53 & 1 & 1 \\
54 & 1 & 1 \\
55 & 1 & 1 \\
56 & 1 & 1 \\
57 & 1 & 1 \\
58 & 1 & 1 \\
59 & 1 & 1 \\
60 & 1 & 1 \\
61 & 1 & 1 \\
62 & 1 & 1 \\
63 & 1 & 3 \\
64 & 2 & 3 \\
65 & 1 & 2 \\
66 & 1 & 2 \\
67 & 2 & 2 \\
68 & 2 & 2 \\
69 & 2 & 2 \\
70 & 2 & 2 \\
71 & 2 & 2 \\
72 & 2 & 2 \\
73 & 2 & 2 \\
74 & 2 & 3 \\
75 & 3 & 3 \\
\end{longtable}
\end{small}

